\subsection{Deleting Names}

The statement \texttt{del name} removes a binding from the current namespace. It does not mean: destroy the object immediately. It means: remove this name's entry from the active namespace dictionary or fast-local slot.

\subsubsection{Removing a Binding with \texttt{del}}

After \texttt{del line}, the name \texttt{line} no longer exists in the current scope. Reading it raises \texttt{NameError: name 'line' is not defined}. The object formerly referenced may still be reachable through other names.

The object and the name are different things. \texttt{del line} deletes the binding named \texttt{line}; it does not directly invoke memory reclamation on the string object.

Compare reassignment to \texttt{None}:

\begin{verbatim}
value = None   # name exists; bound to None singleton
del value      # name no longer exists in this namespace
\end{verbatim}

\subsubsection{Name Deletion, Object Reachability, and Possible Reference Count Changes}

Deleting one name does not necessarily make the object unreachable:

\begin{verbatim}
a = "text"
b = a
del a
# b still references the string object; object remains alive
\end{verbatim}

In CPython, object lifetime is strongly connected to reference counting (\texttt{ob\_refcnt} in \texttt{PyObject\_HEAD}). Removing a binding decreases the count. When the count reaches zero and no cyclic references exist, the object is reclaimed immediately. The programmer should not read \texttt{del name} as a direct \texttt{free()} instruction.

The practical mental model:

\begin{verbatim}
assignment   -> create or change a name-to-object binding
reassignment -> move a name to another object
a = b        -> bind another name to the same object
del name     -> remove one name-to-object binding
\end{verbatim}
